本文基于 PaddleDetection~2.6 框架，在双 Tesla T4 GPU 服务器上对实验鼠目标检测进行了系统研究。
我们从三个异构数据源构建了 7,286 张图像的 VOC 格式数据集（5,828 训练 / 1,458 验证），涵盖两个类别：\textit{mouse}（实验鼠）与 \textit{other}（其他）。
我们设计了 $4\!\times\!2$ 的对照实验矩阵，覆盖两种模型族（PicoDet-S 与 YOLOv3-MobileNetV1）、两种 GPU 配置（单卡 / 双卡）及两种数据规模（全量 / 三分之一子集），共八组训练运行，累计约 270 GPU-小时。

学习率依据线性缩放规则（$\eta = \eta_\text{base}\!\times\! B/B_\text{ref}$）确定，以将批量大小的影响与并行化效应分离。
本文最优模型为 PicoDet-S 单卡小批量训练（batch\,=\,32，lr\,=\,0.08），在 epoch~70 时达到 \textbf{mAP@0.5 = 94.30\%}，较 YOLOv3 最佳结果（91.54\%）高出 \textbf{2.76 个百分点}，同时模型体积 \textbf{缩小 21 倍}（4.4\,MB vs.\ 92\,MB），在相同 GPU 上推理速度 \textbf{提升 2.2 倍}。
一个反直觉的发现是：小批量单卡训练优于标准双卡训练（+0.5\,pp），我们将其归因于梯度噪声在 5,828 样本规模数据集上起到的隐式正则化作用。

训练后的 PicoDet-S 模型被导出为 ONNX 格式，集成至基于 React Native 的应用（ONNX Runtime 与 CoreML 后端），并部署于 iPhone。
经过四轮工程迭代，推理速度由 1\,FPS 提升至稳定 \textbf{14\,FPS}（热机状态）。
我们还研究了数据规模对 YOLOv3 的显著影响：将训练集扩大三倍（1,942~$\to$~5,828 张）使 mAP 提高 \textbf{46.9 个百分点}（44.4\%~$\to$~91.3\%），证实在小领域数据集上，数据工程是锚框式重型架构精度的主要驱动力。
